\section{Local Electric Field}
\label{sec:local}

\subsection{Lorentz Model}

\paragraph{Local field and macroscopic field.}
The response of a material to an external field is complicated, even for a linear isotropic material. The field that polarizes an atom or a molecule, the \textcolor{blue}{local field}, differs from the macroscopic field $\vec E$ of Section~\ref{sec:dielectric}. Besides the external field and the field of the bound charge on the outer surface of the sample, each atom feels the fields of its neighboring dipoles. The macroscopic field is an average over a region that contains many atoms, including the space between them. The local field, in contrast, is the field at the position of one atom, without the field of that atom itself.

\paragraph{Microscopic, macroscopic, and local fields.}
The point dipoles $\vec p_k$ at the positions $\vec r_k$ and the external sources produce the \textcolor{blue}{microscopic field}
\begin{equation*}
    \vec E_{\text{micro}}(\vec r) = \vec E_{\text{ext}}(\vec r) + \sum_k \vec E_{\text{dip}}(\vec r;\vec p_k,\vec r_k),
    \qquad
    \vec E_{\text{dip}}(\vec r;\vec p_k,\vec r_k) = \frac{3(\vec p_k\cdot\hat u_k)\,\hat u_k - \vec p_k}{|\vec r - \vec r_k|^3},
\end{equation*}
where $\hat u_k = (\vec r - \vec r_k)/|\vec r - \vec r_k|$ and $\vec E_{\text{dip}}$ is the point-dipole field of Section~\ref{sec:polarization}. The sum already contains the field of the bound charge on the outer surface, because this charge consists of the uncompensated ends of the dipole chains (Fig.~\ref{fig:boundcharge}). The macroscopic field and the local field at dipole $j$ follow from $\vec E_{\text{micro}}$ by two different operations. The macroscopic field is its average with the weighting function $w$ of Section~\ref{sec:polarization}. The local field is its value at $\vec r_j$ after the field of $\vec p_j$ itself is removed,
\begin{equation}
    \boxed{
    \begin{gathered}
    \vec E(\vec r) = \int w(\vec r - \vec r\,')\,\vec E_{\text{micro}}(\vec r\,')\,d^3r',\\[4pt]
    \vec E_{\text{loc}}(\vec r_j) = \lim_{\vec r \to \vec r_j}\left[\vec E_{\text{micro}}(\vec r) - \vec E_{\text{dip}}(\vec r;\vec p_j,\vec r_j)\right].
    \end{gathered}
    }
\end{equation}
The limit exists, because the subtraction removes the only term that diverges at $\vec r_j$. The difference between the two fields at the position of dipole $j$ is therefore
\begin{equation*}
    \vec E_{\text{loc}}(\vec r_j) - \vec E(\vec r_j)
    = \lim_{\vec r \to \vec r_j}\left[\vec E_{\text{micro}}(\vec r) - \vec E_{\text{dip}}(\vec r;\vec p_j,\vec r_j)\right]
    - \int w(\vec r_j - \vec r\,')\,\vec E_{\text{micro}}(\vec r\,')\,d^3r'.
\end{equation*}
The first term excludes the field of $\vec p_j$ and takes the field at one point. The second term includes the field of $\vec p_j$ and averages over the width $\ell$. Near $\vec r_j$, the fields of the neighbors vary on the scale of the interatomic distance $a \ll \ell$, so the two terms differ.

\paragraph{Macroscopic equations as averaged microscopic equations.}
The equations of Section~\ref{sec:dielectric} contain no explicit averaging, yet they describe macroscopic fields. The reason is that the average with $w$ commutes with the derivatives in the microscopic field equations. A point dipole $\vec p_k$ at $\vec r_k$ has the charge density $-\vec p_k\cdot\vec\nabla\,\delta(\vec r - \vec r_k)$, so the microscopic charge density is
\begin{equation*}
    \rho_{\text{micro}} = \rho_{f,\text{micro}} - \vec\nabla\cdot\vec P_{\text{micro}},
\end{equation*}
where $\rho_{f,\text{micro}}$ is the microscopic density of the free charges. The microscopic field obeys Gauss's law exactly, $\vec\nabla\cdot\vec E_{\text{micro}} = 4\pi\rho_{\text{micro}}$. The average with $w$ is the convolution $(w * f)(\vec r) = \int w(\vec r - \vec r\,')\,f(\vec r\,')\,d^3r'$. Averaging both sides and moving the average through $\vec\nabla$ gives
\begin{equation*}
    \vec\nabla\cdot\left(w * \vec E_{\text{micro}}\right) = 4\pi\left[w * \rho_{f,\text{micro}} - \vec\nabla\cdot\left(w * \vec P_{\text{micro}}\right)\right]
    \quad\Rightarrow\quad
    \vec\nabla\cdot\vec E = 4\pi\left(\rho_f - \vec\nabla\cdot\vec P\right),
\end{equation*}
which is Gauss's law of Section~\ref{sec:dielectric}. The relation $\vec\nabla\times\vec E_{\text{micro}} = 0$ averages to $\vec\nabla\times\vec E = 0$ in the same way. Every solution of Section~\ref{sec:dielectric} is therefore a macroscopic field. For example, the uniform field $\vec E_1$ inside the dielectric sphere is the average of a microscopic field that diverges at every atom. The linear relation $\vec P = \chi\vec E$ also holds only for macroscopic fields, because $\vec P_{\text{micro}}$ is a sum of delta functions and cannot be proportional to $\vec E_{\text{micro}}$.

Lorentz reconcile this by using an imaginary sphere. A sphere is chosen such that it's bigger than the the lattice spacing so verything outside of it can be considered a continuum and small enough compared to the scale that the field is varying so that it feels a single uniform external field. An tom locatd at the center of the sphere now feel 4 compontents
\begin{align}
    E_loc=E_ext + E_dep + E_bound + E_dip
\end{align}
%use underbrace to note them as external field, material surface boundcharge or depolarization field, bound charge of the spherical cavity and field of the diople inside the sphere

The spherical cavity have oposite surface bound charge to that of a sphere to neutralize the charge. So the bound charge of the spherical cavity would be 
\begin{align}
    \sigma_b = - \vec P \cdot \hat r = -P\cos\theta 
\end{align}
% need a figure here 

Because of the azimuthual symmetry, only the azimuthual component of the Electric field survive and it is 
\begin{align}
    E_bz=\int_0^{\pi} \frac{\sigma_b 2\pi\sin\theta R d\theta R \cos\theta}{R^2} = \frac{4\pi P}{3}
\end{align}

If our material is crytallin and is highly symmetric, preferably cubic, there existh a dipole j at rj that cancel the field of dipole i and ri 

So the local field would just be 
\begin{align}
    E_loc = E_xt +\frac{4\pi P}{3}
\end{align}
which is higher than the external field itself. as said